\documentclass[11pt,a4paper]{article}
\usepackage[T1]{fontenc}
\usepackage[utf8]{inputenc}
\usepackage{lmodern,amsmath,amssymb}
\usepackage[margin=25mm]{geometry}
\usepackage[hidelinks]{hyperref}
\hypersetup{pdftitle={One small-field blocking step for $SU(3)$, part 1: the Gaussian fluctuation integral with },pdfauthor={navstokgap research working notes}}
\newcommand{\tightlist}{\setlength{\itemsep}{0pt}\setlength{\parskip}{0pt}}
\setlength{\emergencystretch}{3em}
\setcounter{secnumdepth}{0}
\begin{document}
\hypertarget{one-small-field-blocking-step-for-su3-part-1-the-gaussian-fluctuation-integral-with-block-averaging-with-explicit-constants}{%
\section{\texorpdfstring{One small-field blocking step for \(SU(3)\),
part 1: the Gaussian fluctuation integral with block averaging, with
explicit
constants}{One small-field blocking step for SU(3), part 1: the Gaussian fluctuation integral with block averaging, with explicit constants}}\label{one-small-field-blocking-step-for-su3-part-1-the-gaussian-fluctuation-integral-with-block-averaging-with-explicit-constants}}

This opens the line set in \href{reasons-to-stop-as-research.md}{the
research-directions note}: one Euclidean blocking step of the
small-field renormalization for \(SU(3)\) with every constant written
down. Part 1 is the Gaussian level, where the Wilson action is expanded
to second order and the fluctuation integral is exact. Four things are
established. \textbf{(i)} A block averaging \(W\) of link variables that
intertwines gauge transformations, \(Wd=\bar dW_0\), so that the coarse
field is a gauge field. \textbf{(ii)} In Feynman gauge the quadratic
Wilson form separates into \(4\times8=32\) identical scalar problems,
one per direction and colour, and the fluctuation integral is the
conditional law of a massless free field given its tube averages.
\textbf{(iii)} A block Poincaré inequality with an explicit constant,
\[\langle f,-\Delta f\rangle\ \ge\ \frac49\,\|f\|^2\qquad\text{on tube-mean-zero }f,\]
so the fluctuation covariance is at most \(\tfrac94g^2\) and the
root-mean-square fluctuation of one link angle is at most \(\tfrac32g\);
the alternating modes show the true constant lies in \([\tfrac49,4]\).
\textbf{(iv)} The coarse quadratic form is the Wilson form at the same
\(g^2\) for slowly varying fields, the classical scale invariance of
four dimensions, and its correction and the fluctuation propagator decay
exponentially. For the decay rate the generic Combes--Thomas argument
gives about \(7\times10^{-4}\) per lattice unit against a true rate of
order one: \textbf{the weak side's sharpness problem begins at the
Gaussian propagator}, and the block random-walk expansion is the defined
next piece. Units: \(\hbar\) and \(c\) enter only through \(S_E/\hbar\),
and all lattice quantities are dimensionless with \(a=1\) where not
shown. Constants explicit; nothing promoted.

\hypertarget{conventions}{%
\subsection{1. Conventions}\label{conventions}}

Fine lattice \(\mathbb Z^4\), spacing \(a\); link angles
\(\theta_\ell\in\mathfrak{su}(3)\), \(U_\ell=e^{i\theta_\ell}\), so
\(\theta=aA\) with \(A\) the gauge potential and \(g\) dimensionless.
With \(\operatorname{tr}T^aT^b=\tfrac12\delta^{ab}\),
\[\frac{S_W}{\hbar}=\frac2{g^2}\sum_p\big(N-\operatorname{Re}\operatorname{tr}U_p\big)
=\frac1{2g^2}\sum_p\sum_{a=1}^8\big((d\theta)^a_p\big)^2+O(\theta^3),\]
where
\((d\theta)_p=\theta_{x,\mu}+\theta_{x+\hat\mu,\nu}-\theta_{x+\hat\nu,\mu}-\theta_{x,\nu}\)
is the lattice exterior derivative. The quadratic part is eight copies
of the lattice Maxwell form \(S_2(\theta)=\frac1{2g^2}\|d\theta\|^2\),
and the fluctuation amplitude is of order \(g\): a field strength
\(\eta\) per plaquette costs \(\eta^2/(2g^2)\), in agreement with the
large-field bound \(e^{-\eta^2/(4g^2)}\) of
\href{large-field-action-lower-bound.md}{the lower-bound note} in its
continuum normalization. The cubic and quartic terms, the Haar measure
correction \(dU=(1+O(\theta^2))d\theta\) and the non-abelian part of the
averaging are the remainder, treated in part 2.

\hypertarget{the-averaging-and-its-intertwining-property}{%
\subsection{2. The averaging and its intertwining
property}\label{the-averaging-and-its-intertwining-property}}

Blocks \(B\) of side \(2\); coarse links \((B,\mu)\) from \(B\) to
\(B'=B+2\hat\mu\). Define, for each colour component,
\[(W\theta)_{B,\mu}=\frac1{16}\sum_{x\in B}\big(\theta_{x,\mu}+\theta_{x+\hat\mu,\mu}\big),\]
the average over the sixteen sites of \(B\) of the two-link path from
\(x\) to \(x+2\hat\mu\); and for site functions the block mean
\((W_0\lambda)_B=\frac1{16}\sum_{x\in B}\lambda_x\).

\textbf{Lemma 1 (intertwining).} \(W\,d=\bar d\,W_0\), where \(\bar d\)
is the coarse exterior derivative.

\emph{Proof.}
\((Wd\lambda)_{B,\mu}=\frac1{16}\sum_{x\in B}\big[(\lambda_{x+\hat\mu}-\lambda_x)+(\lambda_{x+2\hat\mu}-\lambda_{x+\hat\mu})\big] =\frac1{16}\sum_{x\in B}(\lambda_{x+2\hat\mu}-\lambda_x)=(W_0\lambda)_{B'}-(W_0\lambda)_B.\)
\(\square\)

So a fine gauge transformation \(\lambda\) moves the coarse field by the
coarse gauge transformation \(W_0\lambda\), and coarse Wilson loops are
functions of fine gauge-invariant data. This is the property that makes
the gauge fixing of Section 3 harmless for the coarse theory.

The tube of \((B,\mu)\) is the set of fine links entering the average:
\(T_{B,\mu}=\{(y,\mu):y\in B\cup(B+\hat\mu)\}\), three \(\mu\)-layers of
eight links with weights \(w=(1,2,1)\), the middle layer counted from
both \(x=y\) and \(x=y-\hat\mu\). Each fine link lies in one or two
tubes.

\hypertarget{feynman-gauge-thirty-two-scalar-problems}{%
\subsection{3. Feynman gauge: thirty-two scalar
problems}\label{feynman-gauge-thirty-two-scalar-problems}}

Add the gauge-fixing term \(\frac1{2g^2}\|d^*\theta\|^2\). Then
\(S_2+S_{\rm gf}=\frac1{2g^2}\sum_\mu\sum_a\langle\theta^a_\mu,-\Delta\theta^a_\mu\rangle\),
because \(d^*d+dd^*\) is the scalar Laplacian on each component of a
lattice one-form. The averaging \(W\) acts on each component separately,
so the Gaussian step is \(32\) copies of one scalar problem: a massless
free field \(f\) on \(\mathbb Z^4\) with covariance
\(g^2(-\Delta)^{-1}\), conditioned on its tube averages
\(\bar f_B=\frac1{16}\sum_{y\in T_B}w_yf_y\) for every block \(B\), with
the tubes elongated along the direction \(\mu\) of the component. The
four directions are equivalent by lattice symmetry.

For gauge-invariant observables the Feynman and Landau measures agree;
for the coarse field they differ by a fine pure gauge \(d\lambda\),
which by Lemma 1 moves the coarse field by the coarse pure gauge
\(\bar dW_0\lambda\), so the law of every coarse gauge-invariant
observable is the same in both. That is what the intertwining buys.

\textbf{Decomposition.} Let \(H\bar f\) be the minimizer of
\(\langle f,-\Delta f\rangle\) subject to \(Wf=\bar f\), and
\(\xi=f-H\bar f\), so \(W\xi=0\). Then, exactly,
\[\langle f,-\Delta f\rangle=\langle H\bar f,-\Delta H\bar f\rangle+\langle\xi,-\Delta\xi\rangle,
\qquad H=GW^T(WGW^T)^{-1},\quad G=(-\Delta)^{-1},\] the cross term
vanishing by the variational characterization. The fluctuation \(\xi\)
is Gaussian on \(K=\{W\xi=0\}\) with covariance
\[C_{\rm fl}=g^2\big[G-GW^T(WGW^T)^{-1}WG\big]
=g^2\big(P(-\Delta)P\big)^{-1}\Big|_K ,\] the conditional covariance of
the free field given its tube averages, \(P\) the orthogonal projection
onto \(K\); and the coarse quadratic form is
\[\bar S_2(\bar f)=\frac1{2g^2}\big\langle\bar f,(WGW^T)^{-1}\bar f\big\rangle
=\min\Big\{\frac1{2g^2}\langle f,-\Delta f\rangle:\ Wf=\bar f\Big\}.\]

\hypertarget{the-block-poincaruxe9-inequality}{%
\subsection{4. The block Poincaré
inequality}\label{the-block-poincaruxe9-inequality}}

\textbf{Proposition 1.} If \(\sum_{y\in T_B}w_yf_y=0\) for every block
\(B\), then
\[\|\partial f\|^2=\langle f,-\Delta f\rangle\ \ge\ \frac49\,\|f\|^2 .\]

\emph{Proof.} One tube \(T\) is a \(2\times2\times2\times3\) box of
\(24\) sites whose graph Laplacian has smallest nonzero eigenvalue
\(\lambda_2=\min(\lambda_2(P_2),\lambda_2(P_3))=\min(2,1)=1\), so for
the plain mean \(m\) of \(f\) on \(T\),
\(\sum_T(f-m)^2\le\sum_{\text{bonds}\subset T}(\partial f)^2\). The
weighted constraint gives \(m=-\frac1{32}\sum_T(w_y-\bar w)(f_y-m)\)
with \(\bar w=\tfrac43\), and
\(\sum_T(w_y-\bar w)^2=8(\tfrac19+\tfrac49+\tfrac19)=\tfrac{16}3\), so
\(|m|\le\frac{4/\sqrt3}{32}\|f-m\|_T\) and
\(24m^2\le\tfrac18\|f-m\|_T^2\). Hence
\(\sum_Tf^2=\|f-m\|_T^2+24m^2\le\tfrac98\sum_{\text{bonds}\subset T}(\partial f)^2\).
Summing over all tubes of the direction \(\mu\): every site lies in at
least one tube, and every bond lies in at most two, so
\(\|f\|^2\le\tfrac98\cdot2\,\|\partial f\|^2\). \(\square\)

\textbf{Consequences.} On \(K\) the precision
\(P(-\Delta)P/g^2\ge\frac4{9g^2}\), so
\[\|C_{\rm fl}\|\le\frac94\,g^2,\qquad
\big\langle\xi_\ell^2\big\rangle^{1/2}=C_{\rm fl}(\ell,\ell)^{1/2}\le\frac32\,g\]
for every fine link and colour. The modes alternating in sign between
consecutive \(\mu\)-layers, or between consecutive sites in a transverse
direction, satisfy the constraint and have
\(\langle f,-\Delta f\rangle=4\|f\|^2\), so the sharp constant lies in
\([\tfrac49,4]\); the value \(\tfrac49\) is what the tube-by-tube
argument gives and is enough for part 2, whose smallness condition is on
\(\tfrac32g\).

\hypertarget{the-coarse-form-and-the-decay-problem}{%
\subsection{5. The coarse form, and the decay
problem}\label{the-coarse-form-and-the-decay-problem}}

\emph{Scale invariance.} For a slowly varying coarse field, \(H\bar f\)
is the slowly varying fine field with the same values, and \(\bar S_2\)
is the fine Maxwell form evaluated on it. In the coarse angle
\(\theta_c=2\bar f\), the natural variable of a link of length \(2a\),
the result is the coarse Wilson quadratic form
\(\frac1{2g^2}\|\bar d\theta_c\|^2\) with the same \(g^2\), up to
corrections of relative order \(\bar k^2a^2\): the Maxwell action is
scale invariant in four dimensions, and the Gaussian step does not run
the coupling. The running appears in part 2, from the dependence of the
fluctuation determinant on the background through the cubic and quartic
vertices.

\emph{Decay.} Both \(C_{\rm fl}\) and the correction to the coarse form
decay exponentially, because the conditioning removes the
long-wavelength modes; the massless pole of \(G\) cancels exactly in
\(G-GW^T(WGW^T)^{-1}WG\) at \(\bar k\to0\). An explicit rate is the
first place where the weak side loses sharpness. The generic
Combes--Thomas argument on
\(\tilde A=P(-\Delta)P+\tfrac49(1-P)\ge\tfrac49\), conjugated by
\(e^{\kappa x\cdot e}\), needs the conjugation error below \(\tfrac29\);
the Laplacian contributes \(2\sinh\kappa\), and the projection, through
\((WW^T)^{-1}=\big(48+8(S+S^{-1})\big)^{-1}\) along \(\mu\), contributes
\(\|P_\kappa-P\|\le10\kappa+O(\kappa^2)\) multiplied by
\(2\|{-\Delta}\|\simeq33\), so
\[\kappa\ \le\ \frac{2/9}{2+330}\ \simeq\ 7\times10^{-4}\ \text{per lattice unit},
\qquad |C_{\rm fl}(x,y)|\le\frac92\,g^2\,e^{-\kappa|x-y|_\infty}.\] The
true rate is of order one per block, as the exact cancellation of the
pole and the alternating-mode energies indicate; the factor of a
thousand is the generic argument's price for the nonlocality of \(P\)
and the size of \(\|{-\Delta}\|\). Balaban's propagator paper (Commun.
Math. Phys. 95 (1984) 17; metadata level) obtains the order-one rate by
a random-walk expansion adapted to the block structure; carrying that
expansion with explicit constants is the next piece, and every
downstream constant, the locality of the effective interaction and the
polymer activities of part 2, inherits its rate.

\hypertarget{what-part-2-needs-from-here}{%
\subsection{6. What part 2 needs from
here}\label{what-part-2-needs-from-here}}

The remainder \(S_W-S_2\) on the small-field region, the Haar correction
and the non-abelian averaging are functions of \(H\bar\theta+\xi\) with
\(\xi\) Gaussian of size \(\tfrac32g\) per link; their expansion is a
polymer gas whose activities carry powers of \(g\) and whose locality
carries the decay rate of \(C_{\rm fl}\). The large-field region
\(|\theta_\ell|\gtrsim1\) has probability at most \(e^{-c/g^2}\) per
link with \(c\) of order \(\tfrac29\) from the covariance bound,
consistent with \href{large-field-action-lower-bound.md}{the lower-bound
note}. The deliverable of part 2 is the Kotecký--Preiss threshold
\(g_{\rm pert}^2\) for this gas.

\hypertarget{consequence-for-state}{%
\subsection{7. Consequence for STATE}\label{consequence-for-state}}

Part 1 of the explicit small-field step is in place: an intertwining
averaging, the reduction to one scalar conditional-Gaussian problem in
Feynman gauge, the block Poincaré constant \(\tfrac49\) with fluctuation
size \(\tfrac32g\) per link, scale invariance of the coarse form, and
the decay problem located with its generic rate \(7\times10^{-4}\)
against a true rate of order one. Next: the block random-walk expansion
for the decay rate, then part 2.
\end{document}
